\section{Discussion}

The compiler shows that a lossless, isotropic dielectric of bounded contrast is
expressive enough to represent every polynomially conditioned sparse linear
system, at the price of a polynomial increase in size. Locality, reciprocity
and a scalar, positive permittivity restrict which couplings a structure can
realize directly, but they do not restrict which equations it can encode once
auxiliary modes and space are available. In this sense, the material
distribution becomes a program: the compiler specifies the structure together
with its source, its readout and the accuracy to which its response must be
known.

For optical computing, the result complements electromagnetic circuit theory
and reconfigurable equation solvers (\cref{tab:prior}).
\citet{MiltonSeppecher2010} establish which terminal responses ideal circuits
can realize, while \citet{tzarouchis2022} give general matrix-inversion
architectures and convergence bounds. Here, the compiler realizes sparse
systems in a continuum dielectric with fixed permeability and bounded scalar
permittivity, and controls errors from the electromagnetic problem through to
the recovered solution. It accounts for the costs that decide whether an
analog solver can be useful as the problem
grows~\citep{Aaronson2015FinePrint,McMahon2023}: programming the structure,
its size, the precision of every coefficient and the readout. The construction
retains the full sparse input class while requiring only logarithmic precision
per coefficient. The same accounting has
been decisive for quantum linear-system
algorithms~\citep{HHL2009,ChildsKothariSomma2017,Aaronson2015FinePrint}, whose
application to electromagnetic scattering and finite-element problems yields at
most polynomial speedups once input and output are
included~\citep{clader2013,MontanaroPallister2016}. Unlike optical Ising
machines, which target NP-hard
optimization~\citep{Inagaki2016,Pierangeli2019,Mohseni2022}, the problem here
is solvable in polynomial time, and the relevant question is the exponent.

\begin{table}[tbp]
\centering
\caption{\textbf{Physical linear-equation solvers and electromagnetic response
synthesis.} The problem column describes the scope of each construction or
architecture; the example column gives the largest reported size $n$ in
experiment (expt.) or simulation (sim.), where available. The table also lists
the physical mechanism -- an iterative feedback loop or a stationary
response -- together with hardware scaling where stated, and the treatment of
accuracy and input and output costs. Theoretical constructions are marked as
such. This work gives a sparse-system reduction with bounds on size, encoding
and decoding cost and accuracy. The numerical illustrations in
\cref{fig:numerical} concern small restricted systems; physical fabrication
remains a separate step.}
\label{tab:prior}
\small
\setlength{\tabcolsep}{3pt}
\renewcommand{\arraystretch}{1.15}
\begin{tabularx}{\linewidth}{@{}>{\raggedright\arraybackslash}p{25mm}>{\raggedright\arraybackslash}X>{\raggedright\arraybackslash}p{19mm}>{\raggedright\arraybackslash}p{32mm}>{\raggedright\arraybackslash}p{32mm}@{}}
\toprule
Work & Problem & Largest example & Mechanism & Accuracy and I/O \\
\midrule
\citet{MiltonSeppecher2010} & Prescribed symmetric terminal response & -- (theory) & Stationary, dielectric and magnetic elements & Ideal material limits; finite-parameter approximation open \\
\citet{estakhri2019} & Integral equation, kernel fixed in the structure & 5 (expt.) & Feedback loop, Neumann series & Measured accuracy \\
\citet{Camacho2021} & Several fixed kernels, nonsymmetric & 5 (expt.) & Feedback loop & Measured accuracy \\
\citet{cordaro2023} & Integral equation, fixed symmetric kernel & 3 (expt.) & Feedback in free space & Fit to simulation \\
\citet{tzarouchis2022} & Programmable complex matrix, generalized inverse & 11 (sim.) & Iteration, $2n^2$ active elements & Convergence time in $\kappa$, $\epsilon$ \\
\citet{chen2024} & Programmable matrix, Richardson iteration & 4 (expt.), 16 by blocking & Iteration, $n^2$ interferometers & Accuracy metric; I/O traffic counted \\
\citet{ye2024} & Fixed diffusion operator on a mesh & $5\times5$ (expt.) & Stationary, one node per grid point & Comparison with simulation \\
\citet{Sun2019PNAS} & Programmable matrix, signs by differences & 4 (expt.), 100 (sim.) & Stationary electronic feedback, $\sim n^2$ cells & Empirical, grows with $\kappa$ \\
\midrule
This work & Any sparse nonsingular system, signed and nonsymmetric & -- (theory) & Stationary response, volume $\Ot(m^{3/2})$ & Residual guarantee; encoding and decoding $\Ot(m^{3/2})$ \\
\bottomrule
\end{tabularx}
\end{table}

For electromagnetic simulation, the result places a limit on what general
solvers can achieve. Fast Maxwell solvers exploit structure beyond locality and
a bounded permittivity, e.g., smooth or piecewise-constant geometry, absorbing
boundaries, moderate quality factors or a small number of distinct
materials~\citep{EngquistYing2011CPAM,Martinsson2019,ShinFan2013}. The
constructed structures have none of these properties. They are closed,
lossless cavities with polynomially many weakly coupled resonant modes, and a
solver that is fast on all of them is, by \cref{eq:transfer}, a fast sparse
linear solver. Inverse design operates in exactly this space of permittivity
distributions~\citep{JensenSigmund2011,Molesky2018}, so an optimizer can, in
principle, produce instances that are as hard to simulate as general sparse
systems. In practice, the difficulty of a given device will depend on how far
it is from this worst case, which the present construction does not quantify.

The exponent $3/2$ coincides with the volume needed to embed a general network
with wires of finite thickness in three
dimensions~\citep{KolmogorovBarzdin1967,GromovGuth2012}. This suggests that the
exponent is natural for any construction that routes couplings through local
paths, but it is not a proof of optimality, since wave-mediated couplings
need not behave like wires. A compiler with nearly linear size would strengthen
the consequence for simulation to $\theta<\gamma$. Conversely, a lower bound on
the volume needed to encode general systems would identify a property of
Maxwell's equations that makes them easier than general sparse linear algebra.

The construction is a mathematical existence proof with deliberately
conservative constants. The modulation of the background is
$2^{-100}$, the proven bounds contain large powers of logarithms, and the fixed
constants are astronomically large -- e.g., the additive constant in the
$O(\log m)$ coefficient length has on the order of $2^{2^{50}}$ bits. As such,
the bounds describe the asymptotic scaling and not the size of any realistic
instance, and the permittivity must be specified to a precision that grows with
the problem size.
The model is an ideal lossless cavity at a single frequency. The theorem
relates the stationary response to the algebraic answer but does not address
the time or energy needed to reach and measure that response, nor fabrication
errors, material dispersion, absorption or detector noise. The result
therefore establishes what a dielectric can represent and what a solver must
handle, not a practical speed or energy advantage for optical computing.
The numerical examples use practical parameter choices: the periodic cavity
illustrates signed encoding after whole-device calibration, while the waveguide
illustrates reusable local assembly in a chain with exact transverse symmetry.
Their scope is the encoding mechanism at small sizes; the general
three-dimensional construction and its asymptotic resource bounds follow from
the analysis. Reusable junctions that preserve
their calibrated response in branching networks would connect the two numerical
settings more closely. Other wave platforms -- acoustic, elastic or quantum --
may admit similar compilers with their own interaction and stability analyses.

In summary, we have shown that sparse linear systems and electromagnetic
fields can be translated into each other at a cost of $\Ot(m^{3/2})$, with
every source of error accounted for. The same construction specifies how
general equations can be written into a passive dielectric and which
algebraic difficulty a general-purpose Maxwell solver inherits, providing a
common, quantitative reference point for computing with light and computing
light.
